% TCC — Fundamentação Teórica
\section{Fundamentação Teórica}\label{sec:fundamentacao}

Esta seção delimita os conceitos necessários para compreender a geração de
testes guiada por contratos e o mecanismo de priorização. A versão definitiva
incorporará as fontes primárias já mapeadas na revisão bibliográfica.

\subsection{APIs REST e Especificação OpenAPI}

Neste trabalho, o documento OpenAPI exerce dois papéis. Primeiro, fornece
estrutura para gerar operações, parâmetros e corpos de requisição. Segundo,
atua como oráculo parcial ao descrever status e schemas esperados nas respostas.
Esse oráculo é incompleto: uma API pode aplicar regras de negócio ausentes do
contrato, e o contrato também pode divergir da implementação sem que exista
falha de segurança.

\textit{[Sketch para orientação: incorporar definições e fontes primárias sobre
REST, OpenAPI e contrato como oráculo parcial.]}

\subsection{Teste Negativo e Fuzzing}

Teste negativo exercita entradas inválidas, limites e combinações adversas para
observar se o sistema rejeita dados de forma controlada. O \textit{fuzzing}
pode ser geracional, quando produz casos a partir de um modelo, ou mutacional,
quando transforma entradas existentes. A ferramenta combina as duas ideias:
gera um valor-base do contrato e aplica operadores tipados. Diferentemente de
fuzzers guiados por cobertura, a observação permanece \textit{black-box}.

\textit{[Sketch para orientação: incorporar fontes primárias sobre teste
negativo, robustez e fuzzing geracional/mutacional.]}

\subsection{Fuzzing de APIs REST Guiado por Contrato}

APIs introduzem desafios de dados estruturados, autenticação, dependências
entre recursos, estado e serialização. OpenAPI permite orientar entradas e
comparar respostas, mas não resolve automaticamente essas dependências nem
garante a completude do oráculo. Por isso, rejeições seguras não documentadas,
aceitações indevidas, erros de servidor e divergências de schema precisam ser
descritos separadamente.

\textit{[Sketch para orientação: consolidar aqui as fontes de fuzzing de APIs e
os limites conhecidos de abordagens exclusivamente black-box.]}

\subsection{Lógica Fuzzy e Prioridade Técnica}

Um sistema Mamdani representa conceitos graduais por funções de pertinência,
combina antecedentes em regras linguísticas, agrega suas ativações e produz um
valor contínuo por desfuzzificação. Essa estrutura favorece uma priorização
interpretável, pois permite preservar entradas, pertinências e regras ativadas.

Neste trabalho, a saída representa prioridade técnica da anomalia. Ela não
confirma vulnerabilidade, explorabilidade, impacto de negócio ou urgência real
de correção. O modelo revisado utiliza contradição comportamental, força da
evidência e exposição de controle de acesso como antecedentes. A detecção
contratual ocorre antes da priorização e constitui uma responsabilidade
conceitualmente distinta.

\textit{[Sketch para orientação: incorporar fundamentos de conjuntos fuzzy,
inferência Mamdani e priorização, distinguindo parâmetros heurísticos de
parâmetros apoiados pela literatura.]}
